% Kapitel 7 -- Von Kuonr\^{a}ts Verlangen und seinen heimlichen Treffen

\section{Von Kuonr\^{a}ts Verlangen und seinen heimlichen Treffen}

\mylettrine{B}{ald} wuchs in \gls{Kuonrat von Steinare}s Brust ein heftigeres Verlangen nach weiblicher Gesellschaft. Denn seit er auf \gls{Hardegg} mit eigenen Augen gesehen hatte, wie \gls{Kuonrat von Hardegg} gegen Frauen auftrat und wie Furcht, Schweigen und gesenkte Blicke sich vor diesem Manne sammelten, hielt er solche Zeichen in seiner Torheit für verborgene Gunst. Was der eine aus Rohheit und Gewalt nahm, das meinte der andere mit stolzem Wesen und hartem Wort erringen zu können. Also gedachte er seiner frühen Verlobung und sprach bei sich, wie mir berichtet ward, dass ihm doch eine Braut zugesprochen sei, die er nicht bloß dem Namen nach besitzen wolle.

So ritt er nach \gls{Raabs an der Thaya}, um \gls{Diemut von Raabs} zu sehen. Sie war damals nicht von jener blendenden Schönheit, von der fahrende Sänger so gern lügen, wohl aber von stiller Art, wohlgefügt und gehorsam in allem, was Haus und Vater von ihr forderten. Eben dies reizte ihn umso mehr, denn \gls{Kuonrat von Steinare} liebte nicht nur das, was ihm gefiel, sondern noch mehr das, was sich ihm entzog.

Als er sie zum ersten Male dort besuchte, fand er sie nicht allein, sondern in Gegenwart einer älteren Zofe und zweier junger Mägde, die an Spindel und Stickrahmen saßen. Darüber verfinsterte sich sogleich sein Sinn. Er hatte gehofft, mit seiner Verlobten allein zu sprechen, sie mit Worten zu bedrängen und ihre Sanftmut für Neigung zu halten. Stattdessen musste er zusehen, wie jede seiner Gebärden von fremden Augen gemessen ward.

Da trat er näher und sprach mit halb gesenkter, doch schon ungeduldiger Stimme: \enquote{\gls{Frouwe} \gls{Diemut von Raabs}, heißt man mich hier als Gast willkommen, oder soll ich vor Mägden Rede stehen wie ein Schulknabe? Schickt sie fort. Ich habe mit Euch allein zu sprechen.}

Die Zofe aber, welche älter war als ihre Herrin und, wie es scheint, auch klüger als der junge Ritter, antwortete vor \gls{Diemut von Raabs}, ehe diese selbst ein Wort fand: \enquote{Herr, die Jungfrau bleibt nicht ohne ehrbares Geleit. So will es ihres Vaters Haus.}

Hierüber geriet der junge Ritter sogleich in Hitze. \enquote{Eures Vaters Haus,} sprach er, \enquote{ist mir wohlbekannt. Doch bin ich nicht irgendein fremder Bursche vom Wege. Bei meinem Ross \gls{Veltloufer}, man schneidet mir hier die \gls{Freiheit} kürzer als einem gebundenen Hund. Weicht von hinnen und lasst mich mit meiner Braut reden.}

Da hob \gls{Diemut von Raabs} endlich den Blick und sagte leise, doch fest: \enquote{Herr, zürnt nicht. Was ehrbar ist, bleibt lieber offen als heimlich. Redet, was Ihr reden wollt, vor denen, die mir zugeordnet sind.}

Solche Worte, die ein besonnener Mann vielleicht als Zucht und Anstand erkannt hätte, galten ihm als Trotz. Er biss die Zähne aufeinander und erwiderte: \enquote{Ihr Weiber nennt jede Fessel ehrbar, sobald sie Euch nur lang genug umgelegt ward.} Dann wandte er sich gegen die Zofe und rief: \enquote{Und du, alte Hexe, du hetzest sie wider mich. Hüte dich, dass ich dich nicht lehre, wem hier Gehorsam gebührt.}

Mit diesen Worten ging er davon, mehr gekränkt als abgewiesen, und sein Stolz litt vielleicht noch schärfer als sein Begehren. Denn \gls{Kuonrat von Steinare} konnte es schlechter ertragen, gehindert zu werden, als entbehrt zu bleiben. Ich zweifle nicht, dass gerade dies den Stachel tiefer setzte.

Einige Tage später kehrte er wieder zurück, diesmal nicht allein, sondern mit dem \gls{Knecht} an seiner Seite. Schon vor dem inneren Hof nahm er ihn beiseite und sprach hastig: \enquote{Hör zu, Knecht. Zieh die schnatternde Gans von der Kammer fort. Rede ihr von einem verlorenen Ring, von einem gestürzten Korb, von was du willst. Nur schaff mir Luft.}

Der \gls{Knecht}, der aus Erfahrung wusste, dass aus seines Herrn Begierden selten Gutes wuchs, antwortete zaghaft: \enquote{Herr, wenn dies kund wird, so fällt die Schuld auf mich.}

Doch \gls{Kuonrat von Steinare} fuhr ihn an: \enquote{Schweig. Deine Schuld ist, zu langsam zu gehorchen. Willst du mich auch noch lehren, was recht sei? Was \gls{Freiheit} bedeutet? Fort jetzt. Ist meine \gls{Freiheit} geringer als das Gemurmel einer Zofe oder das Gekrächz irgendeines Kaplanes? Oder eines Bußpriesters der nur seine keusche Buße im Sinn hat? Nein, du wirst tun, was ich sage, und zwar sofort. Oder willst du, dass ich dich bußen lasse?}

Also tat der \gls{Knecht}, wie ihm geheißen war, und mit einfältiger List lockte er die Zofe auf einen anderen Gang des Hauses. So gewann \gls{Kuonrat von Steinare} endlich jene wenigen Augenblicke allein mit \gls{Diemut von Raabs}, nach denen ihn mehr die reihne Lüsternheit als die Liebe verlangte. Was zwischen den beiden in dieser verborgenen Stunde gesprochen und getan ward, will ich nicht breit erzählen. Einige meinten, sie habe sich anfangs geziert und ihm doch nicht widerstanden. Andere sagten, sie sei mehr aus Furcht vor seinem drängenden Wesen als aus eigenem Verlangen still geblieben. Wieder andere sagen, sie ist eine Frau wie jede andere, die auch verlangen verspürt; vor allem Mädchen ihres alters lassen sich leicht den Kopf verdrehen. Der Herr allein weiß, was in jenem Gemach an diesem Tag vollends geschah.

Nur dies scheint mir wahrscheinlich: \gls{Kuonrat von Steinare} redete auch dort nicht wie ein minniglicher Werber, sondern wie einer, der Besitz für Zuneigung hält. Denn man legte ihm Worte in den Mund wie diese: \enquote{Siehst du nun, \gls{Diemut von Raabs}, wie süß die Stille ist, wenn nicht alle Welt zwischen uns tritt? Du bist mir zugesprochen. Warum sollte ich warten, bis alte Männer und Pfaffen mir erlauben, was doch längst mein ist?}

Und sie soll geantwortet haben: \enquote{Herr, was Gott noch nicht gefügt und die Kirche noch nicht gesegnet hat, das soll der Mensch nicht mit Hast vorwegnehmen.}

Solche Mahnung hätte einen verständigen Mann vielleicht gehemmt. Ihn aber reizte sie nur weiter. Von da an suchte er sie immer wieder auf, bald mit Vorwand, bald mit List, bald mit jener groben Süße, die einfältige Herzen für Stärke halten. Nicht selten sandte er den \gls{Knecht} voraus, damit dieser Türen offen finde, Mägde aufhalte oder Botschaften trage. \gls{Diemut von Raabs} wiederum entzog sich ihm nicht immer, sei es aus Schwäche, aus Neugier oder aus jener gewöhnlichen Hinfälligkeit des Fleisches, die keinem Weibe fern ist. Mehr will ich darüber nicht sagen.

Endlich aber vernahm \gls{Wolfhart von Raabs}, dass sich die Verlobten mehr Freiheit nahmen, als einem ehrbaren Hause zuträglich war. Man sagte, sein Kaplan habe ihm die Gefahr mit ernsten Worten vorgehalten und ihn erinnert, wie leicht Schande schneller wachse als Ruhm. Darauf ließ \gls{Wolfhart von Raabs} seinem künftigen Eidam melden: \enquote{Vor der Hochzeit seht Ihr mein Kind nicht mehr ohne Zeugen. Was recht ist, soll recht bleiben.}

Als \gls{Kuonrat von Steinare} dies vernahm, wurde er so heftig im Gemüt, dass er seinen Grimm sogleich an dem \gls{Knecht} ausließ. Vielleicht gab er ihm die Schuld, weil dieser die Zofen nicht oft genug abgehalten hatte; vielleicht brauchte er auch nur einen Niedrigeren, an dem er seine Schmach austoben konnte. Er packte ihn beim Ärmel und schrie: \enquote{HUND! Du hast alles verdorben. Für Mist und Jauche bist du gut genug, also fahr auch dahin und riech die widerlichen Gase!}

Der \gls{Knecht} beteuerte zitternd seine Unschuld, doch \gls{Kuonrat von Steinare} hörte nicht. Er drängte ihn rückwärts gegen eine offene Senkgrube, wohl in der Absicht, ihm eine abscheuliche Strafe zu bereiten. Allein der Boden war schlüpfrig, und wie es oft geschieht, wenn Gott die Hochfahrt des Menschen mit schneller Hand straft, glitt nicht der Knecht, sondern der Herr selbst aus und stürzte mit einem unerquicklich lauten Geräusch in den Unrat.

So lag er dort, besudelt von Haupt bis Fuß, zum Gespött aller, die davon hörten. Noch einen Monat später, so heißt es, haftete ein übler Gestank an Kleidern, Haar und Bart, und die Leute auf \gls{Hardegg} nahmen dies für eine gerechte Vergeltung. Ja, selbst \gls{Kuonrat von Steinare} konnte sich der Einsicht nicht ganz verschließen, dass ihm hier eine Buße widerfahren war, wie sie seinem ungebändigten Sinn wohl anstand.

Also lehrt diese Begebenheit, dass ungezügeltes Verlangen den Verstand verdunkelt und der Mensch oft eben in jene Grube fällt, die er für einen anderen gegraben hat. Ob \gls{Diemut von Raabs} in jener Zeit mehr verführt oder mehr erschreckt ward, wage ich nicht mit Gewissheit zu entscheiden. Sicher aber ist, dass zwischen ihr und \gls{Kuonrat von Steinare} seit diesen Treffen etwas verdorben war, das auch spätere Worte nicht mehr ganz heilen konnten. Ob aber \gls{Kuonrat von Steinare} aus solcher Schmach wahrhaft eine Lehre zog, daran darf man billig zweifeln; denn mancher nimmt die Buße wohl an seinem Leibe hin, ohne dass sein Sinn darum demütiger werde.
